% Hyper-reference / Hyper-links
\documentclass[12pt]{article}

% ---------- Packages ----------
\usepackage{hyperref}

% \hyperref[label]{text}

\newcommand{\myHyperref}[2]{\hyperref[#1]{#2 \ref*{#1}}} % * was used for discard nested links.

\begin{document}

	\title{Getting Started with \LaTeXe}
	
	\author{A. Author, Professor\\ Some Lab, Some Faculty, University of Iran\\ \texttt{prof-email@email.ac.ir}}
	
	\date{\today}
	
	\maketitle	
	\pagebreak
	
	
	
	
	\section{Introduction} \label{intro}
	{\noindent
	This paper is organized as follow:
	
	\TeX typesetting system is reviewed in \hyperref[what is tex?]{Section ?}; % ? means that in simple \hyperref command can not calculate section number. Instead I've declare a new command for this issue above.
	
	% Advantages of \LaTeX is investigated in section \ref{advantages}; % ** old sentence **
	  Advantages of \LaTeX is investigated in \myHyperref{advantages}{Section}.
	
	Editors software are provided in \myHyperref{editors}{Section} and windows and Linux editors are in \myHyperref{editor.win}{Section} and \myHyperref{editor.linux}{Section}.
	
	And conclusion is in \myHyperref{conclusion}{Section}.
	}
	
	
	
	
	
	\section{What is \TeX?} \label{what is tex?}
		This section is about that what is \TeX?
	
	\section{Advantages of \LaTeX}\label{advantages}
		In this section we discuss about advantages of \LaTeX.

	\section{Editors} \label{editors}
		Editors for \LaTeX is vast and we have so many editors. Also we have some special ones for any type of OS systems.
	
	\subsection{Editors in Windows} \label{editor.win} %\label{section.editor.win}
		This section talk about Windows's editors.
	
	\subsection{Editors in Linux} \label{editor.linux} 
		This section talk about Linux editors.
	
	\section{Conclusion} \label{conclusion}
		Conclusion section is here.

\end{document}